% ======================================================================
% Chapter 2 -- Background
% Target: ~5 pages, <=25% of total. British spelling, present tense,
% active voice, "we"/"one". Every heading opens with a lead-in.
% [VERIFY] = key not yet confirmed in .bib.
% ======================================================================

\chapter{Background}
\label{chap:background}

% TODO: opening motivation paragraph (vision-LANGUAGE model not LLM;
%       "state cues"; BC is the algorithm) + chapter lead-in naming the
%       five threads. (Thesis-wide outline goes in the Introduction.)


% ----------------------------------------------------------------------
\section{Imitation Learning and Behavioural Cloning}
\label{sec:bg-il}
In the absence of a reward signal, an agent can still induce a policy from
expert demonstrations through a supervised learning algorithm, a setting
known as behavioural cloning \cite{Bain1995AFF}. The policy is trained to
reproduce the expert's action at each observed state. A well-documented
weakness of this approach is a compounding-error phenomenon: once the agent
commits an early mistake at deployment, it reaches states the expert
demonstrations never covered, where its predictions degrade further and
errors accumulate over the episode \cite{ross2011}. Formally, if the learned
policy incurs an error probability $\varepsilon$ on states drawn from the
expert distribution, the expected cost of behavioural cloning grows as
$\mathcal{O}(\varepsilon T^2)$ over a horizon of $T$ steps, since each
mistake compounds the mismatch between the states the agent visits and those
the expert demonstrated \cite{ross2011}. This offline-to-online gap is
observed empirically in Minecraft behavioural cloning \cite{kanervisto2020},
and it motivates evaluating the policy through in-environment rollouts
(Section~\ref{sec:rollout}) rather than held-out-frame metrics alone.

This vulnerability is sharper still when the policy conditions on a single
observation: with no temporal context, the agent cannot tell where in a
motion it is, so a momentary perceptual ambiguity can divert it into the
out-of-distribution regime described above. R1V-A operates in exactly this
single-frame regime, which motivates the past-action conditioning and
action-chunking mechanisms introduced in Chapter~\ref{chap:method}.

More recent analysis qualifies how fundamental this gap is. The quadratic
error growth is not intrinsic to behavioural cloning but depends on the
training loss: under a logarithmic loss the dependence on the horizon can be
substantially milder \cite{foster2024horizon}, although this guarantee
weakens again for the continuous action components that policies such as
R1V-A must predict \cite{simchowitz2025pitfalls}. The compounding-error view
therefore remains the appropriate framing here.


% ----------------------------------------------------------------------
\section{Minecraft as a Learning Environment}
\label{sec:bg-minecraft}
Minecraft has emerged as a demanding control benchmark for embodied agents,
owing to its open-ended environment, partial observability, long task
horizons, and large diversity of objectives \cite{fan2022minedojo}. Early
agents achieved competent play through symbolic or scripted interfaces that
abstract away perception and low-level control. A more recent line instead
requires the agent to act from raw pixels through the same
keyboard-and-mouse interface a human uses, a substantially harder setting
that VPT established at scale by behavioural cloning on large quantities of
human gameplay video \cite{baker2022vpt}. R1V-A belongs to this
pixel-and-keypress tradition.

Acting through this interface requires a policy to control the same inputs a
human does: a set of binary key presses together with continuous mouse
movement. VPT formalises this as a factored action space of discrete key
toggles plus a two-axis camera signal, and discretises each camera axis into
mu-law-spaced bins so that the long-tailed distribution of human mouse
movements is represented at higher resolution near zero \cite{baker2022vpt}.
R1V-A adopts this action space and bin convention directly
(Section~\ref{action_space}). Human demonstrations for this setting are
commonly logged in the BASALT contractor format, in which each frame is
paired with the action the player executed, with camera magnitudes already
snapped to the VPT bin centres \cite{shah2021basalt}. The recordings used in
this work follow this convention, so the action logs map onto the policy's
output representation without loss.
% TODO: lifshitz2023steve1 is uncited -- place it (language-for-task-spec
%       example) or drop it from the chapter.


% ----------------------------------------------------------------------
\section{Vision-Language Models and the Frozen-Backbone Paradigm}
\label{sec:bg-vlm}
R1V-A treats a pretrained vision-language model as a fixed feature extractor
and trains only a small head on its representations. This section introduces
the two families of vision-language model the experiments compare, which
differ in how they combine image and text, and then sets out why freezing
such a model is a sound choice.

\subsection{Contrastive and Joint Vision-Language Models}
\label{subsec:bg-vlm-families}
Since we wish to compare the contribution of language to an agent's
capabilities, we examine two structurally different ways of relating images
and text. The two backbones we ablate, CLIP and LLaVA, sit at opposite ends
of a fusion spectrum, and the distinction governs how, and whether, a task
prompt can influence the visual representation.

CLIP \cite{radford2021clip} couples two independent encoders. A text encoder
maps a caption to a single embedding, taken from the activations at the
end-of-sequence position, and a vision encoder (a ViT) maps an image to a
single embedding via its pooled representation. Both are projected into a
shared space. Training is contrastive: for a batch of $N$ image-text pairs,
one forms the $N \times N$ matrix of cosine similarities between every text
and every image embedding, then maximises the $N$ diagonal entries, the true
pairs, against the off-diagonal mismatches. The two modalities therefore
never interact within an example; alignment emerges only through comparison
across the batch, and fusion amounts to a single dot product in the shared
space. A consequence we rely on later is that CLIP's image embedding is a
fixed, language-aligned summary of the frame, computed without reference to
any prompt.

LLaVA \cite{liu2023llava} instead injects vision into a language model. It
reuses a CLIP vision encoder to produce a sequence of patch features,
projects them into the word-embedding space of the language model, and
concatenates these visual tokens with the embedded prompt tokens. The
combined sequence passes through the language model, whose self-attention
layers process the image and text tokens as one sequence, so the two
modalities are jointly contextualised at every layer rather than only at the
output. Fusion here is deep rather than late. The follow-up LLaVA-1.5
\cite{liu2024improvedbaselinesvisualinstruction} keeps this design and
strengthens it, replacing the original single linear projection with a
two-layer multilayer perceptron and broadening the instruction-tuning data
with academic visual-question-answering tasks; this is the backbone R1V-A
uses. This difference in fusion motivates our ablation. In CLIP the prompt
enters only as a parallel embedding, whereas in LLaVA it shares a sequence
with the image inside the model, so the question of whether the prompt
carries usable signal is distinct for the two backbones.

\subsection{Freezing Strong Pretrained Features}
\label{subsec:bg-freeze}
Adapting a large pretrained model to a new task need not mean updating its
weights. Full fine-tuning, the standard approach, carries a documented risk:
when the pretraining and downstream distributions differ, updating all
parameters can distort the pretrained features and degrade performance on
inputs that fall outside the fine-tuning data, the very generalisation a
large backbone is meant to provide
\cite{kumar2022finetuningdistortpretrainedfeatures}. Since R1V-A is evaluated
in rollout worlds it never saw during training, preserving that
out-of-distribution generalisation is a direct concern.

Freezing the backbone and training only a lightweight module on top avoids
this risk and has an established precedent in vision-language modelling.
BLIP-2 connects a frozen image encoder to a frozen language model through a
small trainable bridge, reaching competitive performance while leaving both
large components untouched
\cite{li2023blip2bootstrappinglanguageimagepretraining}. The same principle
underlies R1V-A: the vision-language backbone is held fixed and only the
action head receives gradients.

Freezing is the strict end of a spectrum of parameter-efficient adaptation.
Methods such as LoRA occupy the middle, injecting a small number of trainable
low-rank updates into an otherwise frozen backbone, which adapts the model at
a fraction of the cost of full fine-tuning \cite{hu2022lora}. R1V-A takes the
most conservative position on this spectrum, leaving every backbone parameter
unchanged, which both eliminates the distortion risk entirely and makes the
cached-feature training scheme of Chapter~\ref{chap:method} possible.


% ----------------------------------------------------------------------
\section{Vision-Language-Action Models}
\label{sec:bg-vla}
The frozen backbone and trainable action head together form a
vision-language-action (VLA) model: a system that maps visual and
linguistic input to control output. R1V-A descends from a recent line of
such models, and this section places it within that lineage and against its
closest Minecraft relatives.

The VLA pattern emerged first in robotics, where a vision-language backbone
is repurposed to emit actions rather than text. RT-2 established the
template by training a large vision-language model to output robot commands
as if they were language tokens, transferring web-scale visual and semantic
knowledge into control \cite{brohan2023rt2}. OpenVLA showed the approach
scales as an open recipe at the seven-billion-parameter range comparable to
R1V-A's backbone, adapting the model with low-rank updates rather than full
retraining \cite{kim2024openvla}. Closer to R1V-A's design, RoboFlamingo
keeps the vision-language backbone largely fixed and trains a separate
policy head on top \cite{li2023roboflamingo}, and LLaRA builds a VLA
directly on a LLaVA-style backbone \cite{li2024llara}, while $\pi_0$
attaches a lightweight action module to a frozen model \cite{black2024pi0}.
R1V-A shares this frozen-backbone, separate-head philosophy, but applies it
to Minecraft rather than physical manipulation.

Within Minecraft, two systems frame the design space this thesis works in. PR2L
is the closest published precedent and the methodological template for this
work: it pairs a frozen vision-language model with a learned policy to act from
pixels \cite{chen2024pr2l}, and as in PR2L the contribution here is comparative
and diagnostic --- which frozen backbone yields a usable signal, and how that
signal behaves in closed-loop rollouts --- rather than a new architecture.
JARVIS-VLA sits at the opposite end of the same axis: it fully post-trains the
vision-language backbone on Minecraft data \cite{li2025jarvisvla}, whereas R1V-A
leaves the backbone entirely frozen. It is a conceptual contrast rather than a
benchmarked baseline --- this thesis reproduces no JARVIS-VLA result and does not
measure the full post-training-versus-frozen trade-off that a true head-to-head
would require, confining its conclusions to the frozen regime it actually
studies.

\subsection{Does the Language Pathway Help?}
\label{subsec:bg-language}
A VLA receives a language prompt alongside its visual input, but it does not
follow that the prompt changes what the agent does. Prior work on
language-conditioned imitation finds that a natural-language goal can improve
a policy when tasks must be distinguished, yet image-only policies often
remain strong baselines, since much of the task is already legible from the
visual state alone \cite{lynch2021langlfp}. For R1V-A the question is
sharpened by the frozen-backbone setting: the prompt can only help if the
fixed backbone already encodes it in a form the action head can read.
Whether it does so is what the prompt ablation in Chapter~\ref{chap:method}
measures.

\subsection{Action Chunking}
\label{subsec:bg-chunking}
The compounding-error problem of Section~\ref{sec:bg-il} grows with the
number of independent decisions an agent makes over an episode. Action
chunking attacks this directly: instead of predicting one action per step,
the policy predicts a short sequence of future actions in a single forward
pass, executes the first, and re-plans on the next observation. Introduced by
ACT for fine-grained imitation, this shortens the effective decision horizon
and lets the policy commit to locally coherent motions rather than
re-deciding from scratch at every frame \cite{zhao2023act}. R1V-A adopts the
predict-a-chunk, execute-the-first recipe, which is a particularly natural
fit for its single-frame setting, where consecutive observations carry
little new information and short bursts of action are highly predictable.


% ----------------------------------------------------------------------
\section{Factored Action Spaces and Multi-Label Prediction}
\label{sec:bg-action}
A policy that controls a keyboard faces a combinatorial output space: with
$k$ keys, the set of possible key combinations is $2^k$, far too large to
treat as a single categorical choice. The standard remedy is to factor the
action space, predicting each component independently rather than enumerating
their joint configurations.

Action branching applies this idea by giving each discrete action dimension
its own output head, so the number of predictions grows linearly in the
number of keys rather than exponentially \cite{tavakoli2018actionbranching}.
R1V-A predicts each key as an independent binary decision, which also matches
the structure of the data: a human routinely holds several keys at once, for
example moving forward while attacking, so the targets are genuinely
multi-label rather than mutually exclusive. The head is therefore trained
with a binary cross-entropy loss applied per key, the natural objective for
independent binary outputs.

The two continuous camera axes are handled separately. Following the VPT
convention, each axis is discretised into mu-law-spaced bins and predicted as
a categorical distribution, so that camera control becomes a classification
problem alongside the keys rather than a regression
\cite{baker2022vpt}. This keeps the entire action head within a single
classification framework and lets the mu-law spacing devote finer resolution
to the small movements that dominate human play.


% ----------------------------------------------------------------------
% OPEN ITEMS (still to do):
%  - Write the chapter opening paragraph + chapter lead-in (top of file).
%  - Place or drop lifshitz2023steve1 (2.2).
%  - Add zhao2023act to .bib.
%  - Resolve all [VERIFY] keys against the .bib.
%  - MineDreamer provenance -> Method Dataset section (not here).
%  - Check VLA acronym is not also expanded in the Introduction.
%  - OPTIONAL: justify plain BCE vs reweighted losses (focal/asymmetric)
%    in 2.5 only if you want to cite ridnik2021asymmetric / lin2017focal.
% ----------------------------------------------------------------------